% Part: second-order-logic
% Chapter: syntax-and-semantics
% Section: size-of-domain

\documentclass[../../../include/open-logic-section]{subfiles}

\begin{document}

\olfileid{sol}{syn}{siz}

\olsection{د نامتناهي او \usetoken{S}{enumerable}
  \usetoken{P}{domain} بيانول}

سټ~$M$ هله او يوازې هله (د ډېدېکېند په معنا) نامتناهي دے چې داسې
!!a{injective} تابع~$f\colon M \to M$ وي چې !!{surjective} نۀ وي،
يعنې $\ran{f} \neq M$۔ په لومړۍ درجې منطق کښې يو ځايه
!!{function}~$f$ کتلے شو او ويلے شو چې په
!!a{structure}~$\Struct{M}$ کښې هغې ته ګومارل شوې
تابع~$\Assign{f}{M}$ !!{injective} ده او $\ran{f} \neq
\Domain{M}$:
\[
\lforall[x][\lforall[y][(\eq[f(x)][f(y)] \to \eq[x][y])]] \land
\lexists[y][\lforall[x][\eq/[y][f(x)]]].
\]
که $\Struct{M}$ په دې !!{sentence} کښې صدق وکړي،
$\Assign{f}{M}: \Domain{M} \to \Domain{M}$ !!{injective} ده؛ نو
$\Domain{M}$ بايد نامتناهي وي۔ که $\Domain{M}$ نامتناهي وي، او
داسې تابع وي، $\Assign{f}{M}$ هغه تابع ټاکلے شو، او $\Struct{M}$ به
په !!{sentence} کښې صدق وکړي۔ خو د دې دپاره اړينه ده چې زموږ ژبه
غيرمنطقي سمبول~$f$ ولري چې دې موخې ته يې کاروو۔ په دويمې درجې
منطق کښې په ساده ډول ويلے شو چې داسې تابع \emph{شته}۔ بيا~$f$ ته
اړتيا نۀ وي، او د سوچه دويمې درجې منطق دا !!{sentence} ترلاسه کوو:
\[
\fn{Inf} \ident \lexists[u][(\lforall[x][\lforall[y][(\eq[u(x)][u(y)]
      \lif \eq[x][y])]] \land \lexists[y][\lforall[x][\eq/[y][u(x)]]])].
\]
$\Sat{M}{\fn{Inf}}$ هله او يوازې هله چې $\Domain{M}$ نامتناهي
وي۔ بيا $\fn{Fin} \ident \lnot \fn{Inf}$ تعريفولے شو؛
$\Sat{M}{\fn{Fin}}$ هله او يوازې هله چې $\Domain{M}$ متناهي وي۔
د سوچه لومړۍ درجې منطق يوه !!{sentence} دا نۀ شي بيانولے چې
!!{domain} نامتناهي ده، که څۀ هم د هغو يو نامتناهي سټ يې بيانولے
شي۔ د سوچه لومړۍ درجې منطق د !!{sentence}s هېڅ سټ نشته چې هله او
يوازې هله په !!a{structure} کښې صدق وکړي چې د هغې ساحه متناهي وي۔

\begin{prop}
$\Sat{M}{\fn{Inf}}$ هله او يوازې هله چې $\Domain{M}$ نامتناهي وي۔
\end{prop}

\begin{proof}
$\Sat{M}{\fn{Inf}}$ هله او يوازې هله چې
  $\Sat{M}{\lforall[x][\lforall[y][(\eq[u(x)][u(y)] \lif \eq[x][y])]]
    \land \lexists[y][\lforall[x][\eq/[y][u(x)]]]}[s]$ د کوم~$s$
  دپاره وي۔ که داسې وي، $s(u)$ يوه !!a{injective} تابع ده، او داسې
  $y \in \Domain{M}$ شته چې د~$s(u)$ په مدى کښې نۀ راځي۔ برعکس،
  که داسې !!a{injective} $f\colon \Domain{M} \to \Domain{M}$ وي چې
  $\ran{f} \neq \Domain{M}$، نو $s(u) = f$ د دې ډول متغير ګومارنه
  ده۔
\end{proof}

سټ $M$ هله !!{enumerable} دے چې د !!{element}s يې داسې شمېرنه وي:
\[
m_0, m_1, m_2, \dots
\]
چې تکرار نۀ لري، خو متناهي کېدے شي۔ دا ډول شمېرنه هله او يوازې
هله شته چې يو !!a{element} $z \in M$ او يوه تابع $f\colon M \to M$
وي، داسې چې $z$، $f(z)$، $f(f(z))$، \dots، د~$M$ ټول !!{element}s
وي۔ که شمېرنه وي، $z = m_0$ او $f(m_k) = m_{k+1}$ (يا، که $m_k$
د شمېرنې وروستے !!{element} وي، $f(m_k) = m_k$) اړين !!{element} او تابع
ورکوي۔ بل لور ته، که داسې $z$ او $f$ وي، نو $z$، $f(z)$،
$f(f(z))$، \dots، د~$M$ يوه شمېرنه ده، او $M$ !!{enumerable} دے۔
د $z$ او~$f$ شتون په دويمې درجې منطق کښې بيانولے شو، څو داسې
!!{sentence} جوړه کړو چې هله او يوازې هله په !!a{structure} کښې
رښتيا وي چې !!{structure} !!{enumerable} وي:
\[
\fn{Count} \ident
\lexists[z][\lexists[u][\lforall[X][((X(z) \land
      \lforall[x][(X(x) \lif X(u(x)))]) \lif \lforall[x][X(x)])]]]
\]

\begin{prop}
$\Sat{M}{\fn{Count}}$ هله او يوازې هله چې $\Domain{M}$
!!{enumerable} وي۔
\end{prop}

\begin{proof}
فرض کړئ $\Domain{M}$ !!{enumerable} ده، او $m_0$، $m_1$، \dots،
يې يوه شمېرنه وي۔ د تکرار په لرې کولو ډاډمنېدے شو چې هېڅ $m_k$ دوه
ځلې نۀ راځي۔ $f(m_k) = m_{k+1}$ تعريف کړئ او $s(z) = m_0$ او
$s(u) = f$ وټاکئ۔ ښيو چې
\[
\Sat{M}{\lforall[X][((X(z) \land \lforall[x][(X(x) \lif X(u(x)))])
    \lif \lforall[x][X(x)])]}[s]
\]
فرض کړئ $M \subseteq \Domain{M}$ يو هرومرو فرعي سټ دے۔ بيا فرض
کړئ چې
$\Sat{M}{(X(z) \land \lforall[x][(X(x) \lif
X(u(x)))])}[\Subst{s}{M}{X}]$۔ نو $\Subst{s}{M}{X}(z) \in M$ او
هر کله چې $x \in M$ وي، $(\Subst{s}{M}{X}(u))(x) \in M$ هم وي۔ په
بل عبارت، ځکه چې $\varAssign{\Subst{s}{M}{X}}{s}{X}$، $m_0 \in M$
او که $x \in M$ وي نو $f(x) \in M$؛ نو $m_0 \in M$،
$m_1 = f(m_0) \in M$، $m_2 = f(f(m_0)) \in M$، او داسې نور۔ نو
$M = \Domain{M}$، او په نتيجه کښې
$\Sat{M}{\lforall[x][X(x)]}[\Subst{s}{M}{X}]$۔ ځکه چې
$M \subseteq \Domain{M}$ هرومرو و، ثبوت بشپړ شو:
$\Sat{M}{\fn{Count}}$۔

اوس فرض کړئ $\Sat{M}{\fn{Count}}$، يعنې
\[
\Sat{M}{\lforall[X][((X(z) \land \lforall[x][(X(x) \lif X(u(x)))])
    \lif \lforall[x][X(x)])]}[s]
\]
د کوم~$s$ دپاره۔ $m = s(z)$ او $f = s(u)$ وټاکئ او
$M = \{m, f(m), f(f(m)), \dots\}$ وګورئ۔ داسې تعريف شوے $M$ په
څرګند ډول !!{enumerable} دے۔ نو
\[
\Sat{M}{(X(z) \land \lforall[x][(X(x) \lif X(u(x)))])
    \lif \lforall[x][X(x)]}[\Subst{s}{M}{X}]
\]
د فرض له مخې۔ همدارنګه،
$\Sat{M}{X(z)}[\Subst{s}{M}{X}]$، ځکه چې $M \ni m =
\Subst{s}{M}{X}(z)$، او
$\Sat{M}{\lforall[x][(X(x) \lif
X(u(x)))]}[\Subst{s}{M}{X}]$ هم، ځکه چې که $x \in M$ وي نو
$f(x) \in M$ هم وي۔ نو ځکه چې مقدم او شرطي دواړه صدق کوي، تالي
هم بايد صدق وکړي:
$\Sat{M}{\lforall[x][X(x)]}[\Subst{s}{M}{X}]$۔ خو دا معنا لري چې
$M = \Domain{M}$؛ نو $\Domain{M}$ !!{enumerable} ده، ځکه $M$
د تعريف له مخې همداسې دے۔
\end{proof}

\begin{prob}
!!{sentence} $\fn{Inf} \land \fn{Count}$ په هغو او يوازې هغو
ساحو کښې رښتيا ده چې !!{denumerable} وي۔ د~$\fn{Count}$ تعريف داسې
واړوئ چې يوه بله !!{sentence} شي او په مستقيم ډول وښيي چې ساحه
!!{denumerable} ده، او ثابته يې کړئ۔
\end{prob}

\end{document}
